% =============================================================================
% Section 7 -- Evaluation in Three Live Multi-Agent Environments
%
% Replaces the v8 discrete-event simulation wholesale. All empirical numbers
% enter via the placeholder macros in sections/results_macros.tex; do not
% hand-edit values here. Experiment design authority:
% apip_simulations/apip-testbed-specs-v2.md (+ v2.1 improvements).
% =============================================================================
\section{Evaluation in Three Live Multi-Agent Environments}
\label{sec:simulation}

We evaluate the APIP protocol end-to-end---calibration, trigger, attribution,
tier selection, bounded window, exit---in three live multi-agent environments,
with real LLM inference, controlled failure-mode injection with known ground
truth, and a four-condition baseline family. Each environment was selected to
support a claim the others cannot (Table~\ref{tab:environments}); that the
same contract/trigger/tier/exit machinery governs three unrelated
domains---customer service, software engineering, and economic operations---is
itself the generality argument for the behavioral contract abstraction. An
earlier version of this work evaluated the framework with a parameterized
discrete-event simulation; the present evaluation retires it in favor of
measured behavior.

\begin{table}[t]
  \centering
  \caption{The three evaluation environments and the claim each supports.}
  \label{tab:environments}
  \footnotesize
  \begin{tabular}{@{}L{1.7cm}L{5.6cm}@{}}
    \toprule
    \textbf{Environment} & \textbf{What it buys the evaluation} \\
    \midrule
    $\tau^2$-bench (\texttt{banking\_knowledge})~\cite{barres2025tau2} &
    Customer-service domain match to the deployments that motivate this work;
    deterministic
    database-state ground truth; a retrieval corpus that can be made stale
    independently of task ground truth. \\
    \addlinespace
    ChatDev~\cite{qian2024chatdev} &
    Role-differentiated collaborative mesh (CEO $\to$ CTO $\to$ Programmer
    $\to$ Reviewer $\to$ Tester) with shared context---the substrate for
    alignment creep, disaggregated brittleness detection, and mesh
    disruption. \\
    \addlinespace
    CoffeeBench~\cite{sakana2026coffeebench} \emph{(specified; not run)} &
    Persistent simulated time and a world that evolves without injection:
    \emph{naturally occurring} degradation, which would answer the objection
    that injected drift merely mirrors its own schedule
    (Section~\ref{subsec:eval-coffee}). \\
    \bottomrule
  \end{tabular}
\end{table}

\subsection{Shared Harness and Experimental Design}
\label{subsec:eval-harness}

All three environments plug into a common controller through an adapter
interface; episode traces are normalized to a shared span format (role, phase,
tokens, latency, tool calls and errors, full message content for offline
annotation), so the attribution pipeline is written once and reused
unmodified. Each run proceeds in ticks: a calibration window of clean ticks
derives every contract threshold as mean $\pm\, k\sigma$
(Section~\ref{subsubsec:trigger-stats}); failure modes are then injected on a
pre-registered schedule; both trigger modalities evaluate at every tick (which
fired first is recorded as a result, not short-circuited); breaches pass
through the attribution funnel (Section~\ref{subsubsec:funnel}); and the
selected intervention runs in a bounded window with check-ins at 25/50/75\%,
closing in one of the five exit outcomes of Section~\ref{subsec:exit}.

\textbf{Baseline family.} Every injected cell runs under five conditions:
\emph{APIP} (the full protocol); \emph{B1 no-action} (degradation runs to
horizon); \emph{B2 always-Tier-1} (a scripted prompt patch regardless of
failure mode---the ad-hoc pattern typical of production practice, now one
ablation); \emph{B3 random-tier}; and \emph{B4 oracle} (ground-truth
mode revealed at trigger, correct tier applied immediately). The oracle bounds
achievable performance, making the headline claim positional: APIP achieves
\ttrOraclePct{} of oracle time-to-resolution versus \ttrRandomPct{} for
random-tier selection. Injected ground truth is quarantined from the
controller---it is written into the run record for scoring and revealed only
to B4; a test enforces that the APIP condition never branches on it.

\textbf{Statistical discipline.} Injection schedules, post-calibration
thresholds, the frozen mapping table (hashed into every run manifest), and the
analysis plan are fixed before evaluated runs. Judge models are independent of
backbone models by construction. We report bootstrap confidence intervals
throughout and note that variance across live environments substantially
exceeds what parameterized simulation suggests; we regard reporting that
variance as part of the contribution.

\subsection{$\tau^2$-bench: Behavioral Drift with Deterministic Ground Truth}
\label{subsec:eval-tau}

The primary drift experiment injects nothing into the agent. Instead the
\emph{world} moves: bank policy clauses and the task ground truth advance to
an amended version on a deepening schedule (one additional clause every five
ticks), while the retrieval corpus the agent consults remains frozen at the
old version. The agent's retrieval, reasoning, and tool use all continue to
function; its knowledge source is simply no longer true. This operationalizes
the most familiar production drift incident---a customer-service agent
confidently citing a superseded policy because its knowledge base was never
re-indexed---with the deterministic database providing objective scoring of
every action.

Detection fires after \tauDetectLag{} ticks (judge-scored policy citation
error rate; judge--human agreement \tauCitationJudgeAgree{} on a verified
episode sample). Under APIP, attribution returns \texttt{behavioral\_drift}
and the matched Tier 2 corpus refresh resolves in \tauTTRapip{} ticks. Under
B2, the Tier 1 prompt patch names the amended clauses known at patch time,
recovers briefly, and regresses in \tauBtwoRegressPct{} of runs as the ladder
amends further clauses (TTR \tauTTRbtwo{})---reproducing organically the
patch-then-regress dynamic that makes tier mismatch costly.

\subsubsection{Pilot observations}
\label{subsubsec:eval-tau-pilot}

Before the pre-registered cells, we ran five single-seed pilot iterations of the
drift experiment on local open-weight models (a 7B backbone; an 8B judge from a
different model family), per the gate discipline of not scaling any cell until a
pilot confirms its injector's signature. We report the pilot qualitatively
because its findings shaped the harness and because several of them are evidence
about the \emph{protocol} that scaled runs could not have produced.

\textbf{Detection behaved as designed, and the loop closed durably.} Across
pilot seeds the citation-error trigger fired at lags of 1 to 10 ticks after
injection, absolute-modality first---the modality the taxonomy predicts for
stale-policy drift---and produced zero false breaches across roughly fifty
clean ticks. Attribution returned \texttt{behavioral\_drift} through the
no-split drift rule in every non-degenerate cycle (five matches, five correct
against injected ground truth), and in the final pilot configuration the
matched Tier 2 corpus refresh restored the envelope, with the resolution
holding through the remainder of the horizon. The pilot ground truth then
licensed exactly one adjustment at the mapping-table freeze gate: the no-split
drift rule's prior was raised so that a correct bare-plurality match clears
the abstention threshold; no rule without pilot evidence was touched, and the
frozen table's hash is recorded in every subsequent run manifest.

\textbf{The contract repeatedly caught incomplete remediation.} Our first
Tier~2 implementation updated the experiment's bookkeeping but did not
actually rewrite the documents the agent retrieves---an ecologically familiar
failure: the ticket is closed, the index was never rebuilt. Four ticks after
the ``successful'' remediation, the agent was again citing superseded policy
and the contract re-breached, recording a regression event. Making the refresh
material did not end the lesson: across two further iterations the contract
re-breached on refreshes that were real but \emph{incomplete}---first on
alternate phrasings of the amended clause in other documents, then on a second
document family (the debit-card counterpart of an amended credit-card policy)
that carried the same superseded facts. This is the ad-hoc patch-then-regress
dynamic occurring organically rather than by script, and it surfaces a protocol property we had not designed for
explicitly: because exit evaluation re-measures the behavioral envelope rather
than trusting the intervention log, contract monitoring verifies that a
remediation was \emph{material and complete}, not merely that it was
performed. In Kirkpatrick terms, each incomplete intervention produced a
Level~1/2 record with only partial Level~3/4 change, and the protocol treated
it accordingly. We conjecture that remediation completeness---every phrasing,
every document family, every downstream copy of a superseded fact---is a
dominant failure mode of real-world Tier 2 interventions, and that
contract-based exit evaluation is currently the only systematic check on it.

\textbf{Abstention worked as an honest outcome, twice.} In one pilot iteration
the trace annotator (an 8B judge) systematically mislabeled agent--user
friction as inter-agent misalignment; no mapping rule matched the resulting
label distribution, and the funnel abstained rather than select a tier from
garbage evidence. In another, a post-remediation re-breach attributed at
confidence 0.545 against the 0.55 abstention threshold and was declined. Both
abstentions were logged and scored as attribution failures under our own
accounting---and both were preferable to a confabulated cause.

\textbf{Judge capability is a first-order dependency, and task decomposition
is the lever.} The same 8B judge that consistently failed to detect a verbatim
stale-policy recitation when asked to compare a full transcript against a
policy block answered reliably when the question was decomposed to one atomic
claim against a fact-bearing excerpt (``does the excerpt state that standard
replacement delivery takes 7--10 business days?''). The deployed metric
therefore judges one claim at a time behind a recall-oriented deterministic
pre-filter, with judge misses erring toward false negatives. This finding
directly substantiates the attribution-reliability gap discussed in
Section~\ref{subsec:attribution-scale} and is why the harness enforces
judge--backbone independence and reports judge--human agreement.

\textbf{A design consequence for the reported runs.} An agent that states
\emph{wrong} policy for reasons unrelated to the injection (confabulated
figures, conflated products) contributes a background rate to any judge-scored
error metric. During calibration the citation metric is structurally zero (no
policy has yet diverged), so its baseline cannot absorb that background, and
post-remediation noise is measured against an unrealistically clean null. The
pre-registered runs therefore baseline judge-scored metrics during the
calibration window against the \emph{current} policy version, so that the
trigger tests divergence from the agent's real error floor rather than from
zero.

\subsection{ChatDev: Creep, Brittleness, and Mesh Disruption}
\label{subsec:eval-chatdev}

\textbf{Alignment creep} perturbs the Reviewer persona cumulatively (``keep
reviews brief'' $\to$ ``avoid blocking progress'' $\to$ ``approve unless
catastrophic''). The Reviewer's own flag rate \emph{improves} while downstream
test pass rate decays---the agent whose role is catching failures is the one
degrading, and no absolute floor is crossed. Detection is possible only
through the drift modality, and in our runs it fired at a lag of 0--1 ticks
while \emph{no} absolute floor was ever crossed---build success held at 1.00
throughout---which is the taxonomy's prediction for creep made concrete. The
signature (no dominant trace-level failure mode; falling review engagement)
maps through the frozen table to \texttt{alignment\_creep} and a Tier 1
charter re-injection. In the completed cycle attribution returned
\texttt{alignment\_creep} at confidence 0.73, the restored charter recovered
reviewer engagement to its calibration band within a single tick, and the
recovery held for twelve ticks with no regression (TTR 17).

Two protocol behaviours are worth recording. First, the funnel
\emph{abstained} while only one rung of the ladder had landed: a single step
is genuinely ambiguous evidence, and it committed only once the ramp became
distinguishable---honest abstention that then recovers, rather than abstention
as a dead end. Second, an earlier iteration attributed the degradation to
\texttt{input\_brittleness} and applied Tier 1, which happens to be creep's
matched remediation; the wrong diagnosis administered the right treatment and
erased the injection. Tier correctness and mode correctness are separable, and
only the latter is what the taxonomy claims to deliver.

\textbf{Input brittleness} strips the Programmer's GUI competence for
GUI-category tasks only. Measured against a clean arm, judge-scored
requirement coverage on the targeted category fell from 1.00 to 0.27 while an
untargeted control category was untouched at 1.00---the category-scoped damage
the aggregate-versus-disaggregated comparison depends on.

The consequence is best stated as statistical power rather than a single
detection lag. Required sample size scales as $1/\text{effect}^2$, and a
category holding share $s$ of the batch dilutes its own collapse to $s\times$
the effect in the aggregate; aggregate-only monitoring therefore needs
$1/s^{2}$ times as many episodes to see the same failure. At the 20\% GUI
share of our task mix that is a factor of 25. With our measured effects,
detecting the collapse takes roughly four episodes within the affected group
and roughly one hundred in the aggregate series. Aggregate-only monitoring is
thus not merely slower, as the framing in Section~\ref{sec:taxonomy}
suggests: at realistic batch sizes it is under-powered to see the failure at
all. This also sets the batch size any disaggregated claim requires, since a
group below the contract's power floor receives no envelope
(Section~\ref{subsubsec:trigger-stats}).

\textbf{Mesh disruption} inserts a Documentation Specialist agent between
Programmer and Reviewer at the injection tick. The new agent documents the
Programmer's code and writes the documented version back into the shared code
context, so the Reviewer and Tester operate on a shifted, longer input
distribution. \emph{This experiment returned a null result.} The insertion
executed correctly---the specialist appeared in every post-injection episode,
produced real documentation, and its output propagated into all downstream
phases---yet no peer metric moved: build success held at 1.00, reviewer
engagement stayed inside its calibration band, and tool errors remained at
zero. No contract breach occurred, so no APIP cycle opened and the
scoped-context onboarding protocol could not be evaluated: there was no
disruption to mitigate.

We read this as a precondition finding rather than a refutation. ChatDev
constructs each phase's prompt afresh from templates, so there is no
\emph{bounded} shared context that a new agent can crowd; the
context-crowding mechanism of Section~\ref{subsec:mesh} is structurally
unavailable in this environment. What we could test---distributional shift in
shared state---did not by itself degrade peers. The phenomenon therefore
appears to require an architecture with a bounded or evicting shared context,
which we state as the concrete precondition for future work rather than
leaving the mechanism list untested.

\subsection{Attribution Quality Across All Cells}
\label{subsec:eval-attribution}

Because every injected cell has ground truth, the attribution funnel's output
is scored as a confusion matrix (injected mode $\times$ attributed mode) with
per-mode recall and abstention rate---to our knowledge the first such
accounting for a remediation protocol. Overall accuracy excluding abstentions
is \attrAccuracy{} with an abstention rate of \attrAbstention{}; we report the
two together, never the former alone, since an attribution stage that resolves
every case by construction can achieve any accuracy figure demanded of it.
Only APIP-condition runs enter the matrix: B2 and B3 do not attribute, and
including B4 would render the matrix perfectly diagonal by fiat.

\subsection{CoffeeBench: Designed but Not Run}
\label{subsec:eval-coffee}

CoffeeBench places six agent-operated firms in a persistent economic simulation
with genuine elapsed time, and its role in the design is deliberately narrow
and deliberately irreplaceable: with no injection at all, contract metrics
drift as the world evolves, so the protocol must distinguish degradation that
warrants intervention from legitimate environmental change---the
\texttt{recalibrated} exit outcome of Section~\ref{subsec:exit} exists
because of it, and economic collapse would exercise \texttt{retired}.

We specify this arm but did not run it. A single run costs roughly \$250 in
inference, its 90-day horizon is not truncatable (the documented idle-drift
onset spans days 25--66 and the scheduled demand surge falls on days 40--53),
and substituting cheaper models for all six firms would break comparability
with the published baselines the comparison depends on. We therefore report
the naturally-occurring-degradation question as open. This is the sharpest
remaining limitation of the evaluation: our injected arms cannot by themselves
answer the objection that scheduled degradation mirrors its own schedule.

\subsection{Findings}
\label{subsec:eval-findings}

% Re-derived counterparts of the v8 findings; confirm/refute against real runs.
Pending completion of the pre-registered runs, we structure the analysis
around four questions inherited from the framework's design decisions:
whether tier mismatch remains the dominant source of excess cost (the B2
regression dynamic in Section~\ref{subsec:eval-tau}); whether disaggregation
materially shortens detection (Section~\ref{subsec:eval-chatdev}); whether
token and latency movement leads quality-metric breaches (both are contract
\emph{soft} metrics, logged but not enforced, allowing the leading-indicator
hypothesis to be tested rather than assumed); and whether attribution is
reliable enough to justify the protocol's mandatory attribution phase
(Section~\ref{subsec:eval-attribution}), which is the load-bearing question
for the framework as a whole.

\subsection{Limitations}
\label{subsec:eval-limitations}

Injection schedules in $\tau^2$-bench and ChatDev remain scripted, and the
one environment that would have supplied emergent degradation was not run
(Section~\ref{subsec:eval-coffee}).

A second limitation surfaced during backbone selection and is, we think,
a finding in its own right. Stale-corpus drift is only observable to the extent
that the agent actually restates policy: measuring how often each candidate
backbone asserted any of the amended facts, one model did so in 3.3\% of clean
episodes and another in 37.5\%. On the former, the citation trigger fired once
in twenty-eight post-injection ticks---detection driven by a single episode
rather than a trend. Susceptibility to this failure mode is therefore
model-dependent and inversely related to how readily a model commits to
specifics; a cautious model is both less exposed and harder to monitor. This
strengthens rather than weakens the behavioral-contract argument: which metric
carries signal for a given agent is an empirical question that must be settled
by measurement at deployment time, not assumed. Two contract metrics (policy citation error rate,
requirement coverage) depend on an LLM judge; we report judge--human agreement
on a verified sample and enforce judge/backbone independence, but judge error
remains a bounded unknown. Environment versions are pinned and recorded in
every run manifest---$\tau^2$-bench results in particular are not comparable
across upstream re-grades---and all runs, traces, and the frozen mapping table
are released for replication.
